\documentclass[12pt]{article}
\usepackage{amsmath}
\usepackage{amssymb}
\usepackage{amsfonts}
\usepackage{amsthm}
\usepackage{bm}
\usepackage{array}
\usepackage[margin=1in]{geometry}
\usepackage[colorlinks=true,linkcolor=blue,citecolor=blue,urlcolor=blue]{hyperref}

\numberwithin{equation}{section}

\newtheorem{theorem}{Theorem}
\newtheorem{proposition}{Proposition}
\newtheorem{lemma}{Lemma}
\newtheorem{corollary}{Corollary}
\newtheorem{assumption}{Assumption}
\theoremstyle{definition}
\newtheorem{definition}{Definition}
\newtheorem{remark}{Remark}

\newcommand{\ind}{\mathbf 1}
\newcommand{\R}{\mathbb R}
\newcommand{\E}{\mathbb E}
\newcommand{\Pp}{\mathbb P}
\newcommand{\F}{\mathcal F}
\newcommand{\op}{o_p}
\newcommand{\Op}{O_p}

\begin{document}

\begin{center}
{\Large\bf Primitive Theory Workbook: Solved and Audited Edition}\\[6pt]
{\large Profile Sieve Empirical Likelihood with Nuisance Breaks}\\[6pt]
{\normalsize Complete proofs, exact rate calculations, and necessary repairs}
\end{center}

\tableofcontents
\newpage

\section{Primitive Setup and Notation}

For a scalar array \(\bm a=(a_1,\ldots,a_T)'\), write
\[
    \|\bm a\|_T^2=T^{-1}\sum_{t=1}^T a_t^2,
    \qquad
    \|\bm a\|_\infty=\max_{t\leq T}|a_t|.
\]
All \(O_p(\cdot)\) and \(o_p(\cdot)\) statements for vectors and matrices use Euclidean and spectral norm unless another norm is displayed.

The maintained null model is
\begin{equation}
    Y_t
    =
    \beta_0X_{t-1}
    +
    \sum_{j=0}^{q_0}
    m_j(W_{t-1})\ind\{k_{j,0}<t\leq k_{j+1,0}\}
    +
    U_t,
    \label{eq:primitive-model}
\end{equation}
where
\[
    0=k_{0,0}<k_{1,0}<\cdots<k_{q_0,0}<k_{q_0+1,0}=T.
\]
The coefficient \(\beta_0\) is stable. All structural breaks are in the nuisance component.

For a candidate partition \(\Lambda=(k_1,\ldots,k_q)\), let \(k_0=0\), \(k_{q+1}=T\), and
\[
    I_j(\Lambda)=\{t:k_j<t\leq k_{j+1}\},
    \qquad
    j_\Lambda(t)=j\ \text{if }t\in I_j(\Lambda).
\]
Let
\[
    \mu_t=m_{j_{\Lambda_0}(t)}(W_{t-1}),
    \qquad
    \bm\mu=(\mu_1,\ldots,\mu_T)'.
\]

Let \(\bm P\) be the \(T\times K\) sieve matrix with \(t\)-th row
\[
    \bm P_K(W_{t-1})'=(p_1(W_{t-1}),\ldots,p_K(W_{t-1})).
\]
For a partition \(\Lambda\), define
\begin{equation}
    \bm Q_{K,\Lambda}
    =
    [\bm D_0(\Lambda)\bm P,\ldots,\bm D_q(\Lambda)\bm P],
    \qquad
    \bm M_{K,\Lambda}
    =
    \bm I_T-\bm Q_{K,\Lambda}
    (\bm Q_{K,\Lambda}'\bm Q_{K,\Lambda})^\dagger
    \bm Q_{K,\Lambda}'.
    \label{eq:primitive-qm}
\end{equation}
The Moore--Penrose inverse makes \(\bm M_{K,\Lambda}\) defined for every candidate partition.

For each \((\beta,\Lambda)\), define
\begin{align}
    \widehat{\bm u}(\beta,\Lambda)
    &=
    \bm M_{K,\Lambda}(\bm Y-\beta\bm X_{lag}),
    \\
    \bm w^c(\Lambda)
    &=
    \bm M_{K,\Lambda}\bm w,
    \\
    \widehat Z_t(\beta,\Lambda)
    &=
    \widehat u_t(\beta,\Lambda)w_t^c(\Lambda).
\end{align}
Let
\[
    S_T(\Lambda)
    =
    T^{-1/2}\sum_{t=1}^T\widehat Z_t(\beta_0,\Lambda),
    \qquad
    V_T(\Lambda)
    =
    T^{-1}\sum_{t=1}^T\widehat Z_t(\beta_0,\Lambda)^2.
\]

For a scalar score array \(Z_1,\ldots,Z_T\) satisfying the convex-hull
condition, let \(\lambda_T\) be the unique solution of
\[
    \sum_{t=1}^T\frac{Z_t}{1+\lambda_T Z_t}=0,
    \qquad
    1+\lambda_T Z_t>0\quad(t\leq T),
\]
and define the twice negative empirical-likelihood log ratio by
\begin{equation}
    \ell_T(Z_1,\ldots,Z_T)
    =2\sum_{t=1}^T\log(1+\lambda_TZ_t).
    \label{eq:sol-el-definition}
\end{equation}
In particular,
\[
    \ell_T(\beta,\Lambda)
    =\ell_T\{\widehat Z_1(\beta,\Lambda),\ldots,
                    \widehat Z_T(\beta,\Lambda)\}.
\]

For an interval \([s,e]\), define the null-imposed segment cost
\[
    \mathcal C_K(s,e;\beta)
    =
    \min_{\bm c\in\R^K}
    \sum_{t=s}^e
    \{Y_t-\beta X_{t-1}-\bm P_K(W_{t-1})'\bm c\}^2.
\]
For \(\Lambda=(k_1,\ldots,k_q)\), define
\[
    RSS_K(\beta,\Lambda)
    =
    \sum_{j=0}^q
    \mathcal C_K(k_j+1,k_{j+1};\beta).
\]
All minimizers are measurable selections with a fixed deterministic tie-breaking rule. For known order,
\[
    \widehat\Lambda_q(\beta)
    =
    \arg\min_{\Lambda\in\mathcal L_q(\epsilon)}
    RSS_K(\beta,\Lambda),
    \qquad
    \widetilde\Lambda=\widehat\Lambda_{q_0}(\beta_0).
\]
For unknown order,
\[
    \widehat q(\beta)
    =
    \arg\min_{0\leq q\leq \bar q}
    \{RSS_K(\beta,\widehat\Lambda_q(\beta))+\operatorname{pen}_T(q,K)\},
    \qquad
    \widehat\Lambda(\beta)=\widehat\Lambda_{\widehat q(\beta)}(\beta).
\]

\section{Primitive Assumptions}

\begin{assumption}[Primitive stochastic environment]
\label{ass:prim-stoch}
The triangular array is defined on a filtered probability space
\((\Omega,\mathcal A,\{\F_t\}_{t=0}^T,\Pp)\). The variables \(X_{t-1}\), \(W_{t-1}\), \(\bm P_K(W_{t-1})\), and \(w_t=w(X_{t-1})\) are \(\F_{t-1}\)-measurable. The error satisfies
\[
    \E(U_t\mid \F_{t-1})=0,
    \qquad
    \E(|U_t|^{2+\delta}\mid \F_{t-1})\leq C
\]
for some \(\delta>0\) and finite \(C\). The conditional variance \(\sigma_t^2=\E(U_t^2\mid\F_{t-1})\) is bounded above and away from zero.
\end{assumption}

\begin{assumption}[Score admissibility under full-sample projection]
\label{ass:prim-score-admissible}
For the true partition, \(\bm v_0=\bm M_{K,\Lambda_0}\bm w\) satisfies one of the following primitive routes.
\[
    \text{Route A: } \bm v_0 \text{ is admissible in the martingale CLT below.}
\]
If the full-sample projection is not predictable, assume there is an
\(\F_{t-1}\)-measurable approximation \(v_{t,0}^\ast\) such that
\[
    T^{-1/2}\sum_{t=1}^T U_t(v_{t,0}-v_{t,0}^\ast)=o_p(1),
    \qquad
    T^{-1}\sum_{t=1}^T U_t^2\{v_{t,0}^2-(v_{t,0}^\ast)^2\}=o_p(1).
\]
\end{assumption}

\begin{assumption}[Bounded weights, nuisance functions, and basis envelope]
\label{ass:prim-bounded}
There are constants \(C_w,C_m<\infty\) and an envelope \(\zeta_K\) such that
\[
    \max_{t\leq T}|w_t|\leq C_w,
    \qquad
    \max_{0\leq j\leq q_0}\sup_x|m_j(x)|\leq C_m,
    \qquad
    \max_{t\leq T}\|\bm P_K(W_{t-1})\|\leq \zeta_K
\]
with probability tending to one.
\end{assumption}

\begin{assumption}[Spacing, smoothness, and break detectability]
\label{ass:prim-breaks}
The number \(q_0\) is fixed. For some \(\epsilon>0\),
\[
    \min_{0\leq j\leq q_0}(k_{j+1,0}-k_{j,0})\geq \epsilon T.
\]
Each \(m_j\) belongs to a common smoothness class \(\mathcal H^s\). If \(q_0\geq1\), then
\[
    \Delta_T=\min_{0\leq j<q_0}\|m_{j+1}-m_j\|_{L^2(P_W)}>0,
    \qquad
    \frac{T\Delta_T^2}{K\log T}\to\infty.
\]
\end{assumption}

\begin{assumption}[Uniform Gram stability]
\label{ass:prim-gram}
For the true partition and all candidate partitions in the local neighborhood
\[
    \mathcal N_T(r_T)
    =
    \{\Lambda:\max_j|k_j-k_{j,0}|\leq C r_T,\ q=q_0\},
\]
the eigenvalues of \(T^{-1}\bm Q_{K,\Lambda}'\bm Q_{K,\Lambda}\), restricted to nonzero directions, are bounded away from zero and infinity with probability tending to one.
\end{assumption}

\begin{assumption}[Sieve approximation and projection bias]
\label{ass:prim-sieve}
For each regime \(j\), there is a coefficient vector \(\bm c_{j,K}\) such that
\[
    m_j(W_{t-1})
    =
    \bm P_K(W_{t-1})'\bm c_{j,K}+r_{j,K,t},
    \qquad
    \max_j\|\bm r_{j,K}\|_T=O_p(a_K),
    \qquad
    \max_j\|\bm r_{j,K}\|_\infty=O_p(a_K).
\]
Let \(\bm r_0\) be the stacked true-partition approximation error. The rates satisfy
\[
    T^{-1/2}\bm r_0'\bm M_{K,\Lambda_0}\bm w=o_p(1),
    \qquad
    a_K(1+\zeta_K)\to0,
    \qquad
    a_K(1+\zeta_K)^2=o(\sqrt T).
\]
Also,
\[
    \frac{(q_0+1)K}{\sqrt T}\to0,
    \qquad
    \zeta_K\sqrt{\frac{(q_0+1)K}{T}}\to0.
\]
\end{assumption}

\begin{assumption}[Fixed-order localization target]
\label{ass:prim-localization-target}
For known \(q_0\),
\[
    \max_{1\leq j\leq q_0}|\widetilde k_j-k_{j,0}|=O_p(r_T),
    \qquad
    r_T=o(\sqrt T).
\]
This assumption is a target to be derived in Section \ref{sec:break-localization}; it may be used as an input in Section \ref{sec:one-break-bridge}.
\end{assumption}

\begin{assumption}[Empirical-likelihood convex hull]
\label{ass:prim-convex}
For the score array under consideration,
\[
    \Pp\{\min_{t\leq T}Z_{t,T}<0<\max_{t\leq T}Z_{t,T}\}\to1.
\]
When used for the true partition, \(Z_{t,T}=\widehat Z_t(\beta_0,\Lambda_0)\). When used for the estimated partition, \(Z_{t,T}=\widehat Z_t(\beta_0,\widetilde\Lambda)\).
\end{assumption}


\section{Rigour Audit and Completing Conditions}
\label{sec:rigour-audit}

Several displays in the original workbook were labelled ``targets.''  They
are not all consequences of Assumptions \ref{ass:prim-stoch}--\ref{ass:prim-convex}
as originally written.  The precise gaps are: the predictable approximation
needs an envelope and variance-replacement condition; local partition
perturbations need local stability of the sieve bias; logarithmic RSS
concentration does not follow from only a finite \((2+\delta)\)-moment; and an
overfitting penalty that merely diverges need not dominate the searched RSS
gain.  The following conditions make every proof below logically complete.
They are stated separately so that none of the missing hypotheses is hidden.

Write
\[
    d_K=(q_0+1)K,
    \qquad b_K=1+\zeta_K,
    \qquad \delta_0=\min\{\delta,2\}.
\]

\begin{assumption}[Oracle envelope completion]
\label{ass:sol-oracle-envelope}
When the predictable-approximation route in Assumption
\ref{ass:prim-score-admissible} is used, the approximation can be chosen so
that
\[
    \max_{t\leq T}|v_{t,0}^\ast|=O_p(b_K),
    \qquad
    T^{-1}\sum_{t=1}^T\sigma_t^2
       \{(v_{t,0}^\ast)^2-v_{t,0}^2\}=o_p(1).
\]
Moreover,
\begin{equation}
    b_K^{2+\delta_0}T^{-\delta_0/2}\longrightarrow0.
    \label{eq:sol-envelope-rate}
\end{equation}
If \(\bm v_0\) itself is predictable, take \(v_{t,0}^\ast=v_{t,0}\).
\end{assumption}

\begin{assumption}[Local sieve-bias stability]
\label{ass:sol-local-bias}
In the one-break section, for every fixed \(C<\infty\),
\begin{equation}
    \sup_{|k-k_0|\leq Cr_T}
    T^{-1/2}\left|
      \bm r_0'\{\bm M_{K,k}-\bm M_{K,k_0}\}\bm w
    \right|=o_p(1).
    \label{eq:sol-local-bias}
\end{equation}
A sufficient primitive rate, proved in the solution of Lemma
\ref{lem:wb-bridge-mu}, is
\begin{equation}
    \frac{a_K b_K(1+\zeta_K^2)r_T}{\sqrt T}\longrightarrow0.
    \label{eq:sol-local-bias-sufficient}
\end{equation}
\end{assumption}

\begin{assumption}[Local full-rank series regularity]
\label{ass:sol-local-series}
In the one-break section, every segment Gram matrix associated with a
partition satisfying \(|k-k_0|\leq Cr_T\) is nonsingular, and its eigenvalues
are between \(cT\) and \(CT\), uniformly with probability tending to one.
Also
\begin{equation}
    \zeta_K^2=O(K).
    \label{eq:sol-regular-envelope}
\end{equation}
The last display is the usual normalized-series envelope condition.  It is
needed for the denominator rate stated in Lemma \ref{lem:wb-bridge-v}; without
it that rate must be replaced by the larger order
\(b_K^2\zeta_K^2r_T/T\).
\end{assumption}

\begin{assumption}[Localization concentration completion]
\label{ass:sol-localization}
This assumption is used only in Section \ref{sec:break-localization}.  Given
the design sigma-field \(\mathcal D_T\), the errors are independent, centered,
conditionally sub-Gaussian with a common conditional variance \(\sigma^2\).
Every minimum-spaced one-break sieve matrix has rank \(2K\).  If
\(h=|k-k_0|\) and
\(\bm A_k=\bm M_{K,k}-\bm M_{K,k_0}\), then, uniformly over candidates,
\begin{align}
    \bm\mu'\bm A_k\bm\mu
    &\geq c_\Delta h\Delta_T^2-\rho_T(k),
    &
    \sup_{k:\,h\geq1}\frac{|\rho_T(k)|}{h\Delta_T^2}&=o_p(1),
    \label{eq:sol-profile-identification}
    \\
    \|\bm A_k\bm\mu\|^2&\leq C hK
    \label{eq:sol-mean-sensitivity}
\end{align}
with probability tending to one.  Condition
\eqref{eq:sol-profile-identification} is a sample profile-identification
condition.  It follows, for example, in the exact-sieve case from uniform
segment Gram bounds and a local empirical jump-energy lower bound; this is
verified explicitly in the proof of Lemma \ref{lem:wb-rss-drift}.
\end{assumption}

\begin{remark}[Why the localization completion is necessary]
The martingale restriction
\(\E(U_t\mid\mathcal F_{t-1})=0\) does not imply
\(\E(U_t\mid\mathcal D_T)=0\), because \(\mathcal D_T\) contains future
regressors and those regressors may reveal \(U_t\).  Also, a finite
\((2+\delta)\)-moment does not by itself yield a \(\log T\) search bound for a
quadratic form.  Thus Lemmas \ref{lem:wb-rss-drift} and
\ref{lem:wb-rss-fluctuation} genuinely require an additional route such as
Assumption \ref{ass:sol-localization}.
\end{remark}

\begin{lemma}[Scalar empirical-likelihood expansion]
\label{lem:sol-scalar-el}
Let \(Z_{1,T},\ldots,Z_{T,T}\) satisfy the convex-hull condition and suppose
\[
    A_T=T^{-1/2}\sum_{t=1}^TZ_{t,T}=O_p(1),
    \qquad
    B_T=T^{-1}\sum_{t=1}^TZ_{t,T}^2\geq c+o_p(1)
\]
for some \(c>0\), and
\[
    \max_{t\leq T}|Z_{t,T}|=o_p(\sqrt T).
\]
Then
\begin{equation}
    \ell_T(Z_{1,T},\ldots,Z_{T,T})
    =\frac{A_T^2}{B_T}+o_p(1).
    \label{eq:sol-scalar-el-expansion}
\end{equation}
\end{lemma}

\begin{proof}
Put
\(S_1=\sum_tZ_{t,T}\), \(S_2=\sum_tZ_{t,T}^2\), and
\(M_T=\max_t|Z_{t,T}|\).  The convex-hull condition gives a unique multiplier
\(\lambda_T\) in the interval on which every \(1+\lambda_TZ_{t,T}\) is
positive.  Its score equation implies the exact identity
\[
    S_1
    =\lambda_T\sum_{t=1}^T
      \frac{Z_{t,T}^2}{1+\lambda_TZ_{t,T}}.
\]
Consequently,
\[
    |S_1|
    \geq
    \frac{|\lambda_T|S_2}{1+|\lambda_T|M_T}.
\]
After rearrangement,
\[
    |\lambda_T|\{S_2-|S_1|M_T\}\leq |S_1|.
\]
Now \(|S_1|=O_p(\sqrt T)\), \(S_2\geq(c+o_p(1))T\), and
\(|S_1|M_T=o_p(T)\).  Hence
\[
    \lambda_T=O_p(T^{-1/2}),
    \qquad
    |\lambda_T|M_T=o_p(1).
\]
Using \((1+x)^{-1}=1-x+x^2/(1+x)\) in the score equation gives
\[
    0=S_1-\lambda_TS_2+R_{1T},
    \qquad
    |R_{1T}|
    \leq
    \frac{\lambda_T^2M_TS_2}{1-|\lambda_T|M_T}.
\]
Thus
\[
    \lambda_T=\frac{S_1}{S_2}+o_p(T^{-1/2}).
\]
Finally, Taylor expansion of the logarithm, uniformly because
\(|\lambda_T|M_T=o_p(1)\), yields
\[
    2\sum_t\log(1+\lambda_TZ_{t,T})
    =2\lambda_TS_1-\lambda_T^2S_2+R_{2T},
\]
where
\[
    |R_{2T}|
    \leq C|\lambda_T|^3M_TS_2=o_p(1).
\]
Substituting the expansion for \(\lambda_T\) gives
\[
    2\sum_t\log(1+\lambda_TZ_{t,T})
    =\frac{S_1^2}{S_2}+o_p(1)
    =\frac{A_T^2}{B_T}+o_p(1),
\]
which proves \eqref{eq:sol-scalar-el-expansion}.
\end{proof}

\section{Deterministic Projection Algebra}

\begin{lemma}[Block projection identities]
\label{lem:wb-projection}
For every partition \(\Lambda\),
\[
    \bm Q_{K,\Lambda}'\bm M_{K,\Lambda}=\bm0,
    \qquad
    \bm M_{K,\Lambda}'=\bm M_{K,\Lambda},
    \qquad
    \bm M_{K,\Lambda}^2=\bm M_{K,\Lambda}.
\]
Consequently, for every conformable \(\bm c\),
\[
    (\bm Q_{K,\Lambda}\bm c)'\bm M_{K,\Lambda}\bm w=0.
\]
\end{lemma}


\begin{proof}[Detailed solution]
Abbreviate \(\bm Q=\bm Q_{K,\Lambda}\).  Take a thin singular-value
decomposition
\[
    \bm Q=\bm U_r\bm\Sigma_r\bm V_r',
\]
where \(r=\operatorname{rank}(\bm Q)\), the columns of \(\bm U_r\) and
\(\bm V_r\) are orthonormal, and \(\bm\Sigma_r\) is diagonal with strictly
positive entries.  Then
\[
    \bm Q'\bm Q
    =\bm V_r\bm\Sigma_r^2\bm V_r',
    \qquad
    (\bm Q'\bm Q)^\dagger
    =\bm V_r\bm\Sigma_r^{-2}\bm V_r'.
\]
Therefore
\[
    \bm\Pi_{K,\Lambda}
    =\bm Q(\bm Q'\bm Q)^\dagger\bm Q'
    =\bm U_r\bm U_r'.
\]
This is the orthogonal projector onto \(\operatorname{col}(\bm Q)\).  In
particular,
\[
    \bm\Pi_{K,\Lambda}'=\bm\Pi_{K,\Lambda},
    \qquad
    \bm\Pi_{K,\Lambda}^2=\bm\Pi_{K,\Lambda},
    \qquad
    \bm\Pi_{K,\Lambda}\bm Q=\bm Q.
\]
Since \(\bm M_{K,\Lambda}=\bm I_T-\bm\Pi_{K,\Lambda}\),
\[
    \bm M_{K,\Lambda}'=\bm M_{K,\Lambda}
\]
and
\[
    \bm M_{K,\Lambda}^2
    =(\bm I_T-\bm\Pi_{K,\Lambda})^2
    =\bm I_T-2\bm\Pi_{K,\Lambda}+\bm\Pi_{K,\Lambda}^2
    =\bm M_{K,\Lambda}.
\]
Also
\[
    \bm Q'\bm M_{K,\Lambda}
    =\{\bm M_{K,\Lambda}\bm Q\}'
    =\bm0'.
\]
For conformable \(\bm c\) and \(\bm w\), it follows immediately that
\[
    (\bm Q\bm c)'\bm M_{K,\Lambda}\bm w
    =\bm c'\bm Q'\bm M_{K,\Lambda}\bm w=0.
\]
\end{proof}


\begin{lemma}[Uniform fitted-value and leverage bounds]
\label{lem:wb-leverage}
Under Assumptions \ref{ass:prim-bounded} and \ref{ass:prim-gram}, uniformly over \(\Lambda\in\mathcal N_T(r_T)\),
\[
    \|\bm M_{K,\Lambda}\bm w\|_\infty=O_p(1+\zeta_K),
    \qquad
    \max_{t\leq T}\{\bm Q_{K,\Lambda}(\bm Q_{K,\Lambda}'\bm Q_{K,\Lambda})^\dagger\bm Q_{K,\Lambda}'\}_{tt}
    =
    O_p(\zeta_K^2K/T).
\]
\end{lemma}


\begin{proof}[Detailed solution]
Write \(\bm Q=\bm Q_{K,\Lambda}\), \(\bm G=\bm Q'\bm Q\), and let
\(\bm q_t'\) denote row \(t\) of \(\bm Q\).  Because exactly one regime block
is active in each row,
\[
    \|\bm q_t\|=\|\bm P_K(W_{t-1})\|\leq\zeta_K.
\]
On the uniform Gram event in Assumption \ref{ass:prim-gram},
\[
    \|\bm G^\dagger\|\leq \frac{C}{T},
    \qquad
    \|\bm Q\|^2=\lambda_{\max}(\bm Q'\bm Q)\leq CT.
\]
Since \(\|\bm w\|\leq C_w\sqrt T\), the least-squares coefficient vector
\[
    \widehat{\bm c}_\Lambda=\bm G^\dagger\bm Q'\bm w
\]
satisfies
\[
    \|\widehat{\bm c}_\Lambda\|
    \leq \frac{C}{T}\,\|\bm Q\|\,\|\bm w\|
    \leq C.
\]
Thus, uniformly in \(t\) and \(\Lambda\),
\[
    |(\bm\Pi_{K,\Lambda}\bm w)_t|
    =|\bm q_t'\widehat{\bm c}_\Lambda|
    \leq C\zeta_K,
\]
and hence
\[
    \|\bm M_{K,\Lambda}\bm w\|_\infty
    \leq \|\bm w\|_\infty+
          \|\bm\Pi_{K,\Lambda}\bm w\|_\infty
    =O_p(1+\zeta_K).
\]
For the leverage,
\[
    (\bm\Pi_{K,\Lambda})_{tt}
    =\bm q_t'\bm G^\dagger\bm q_t
    \leq \|\bm G^\dagger\|\,\|\bm q_t\|^2
    \leq C\frac{\zeta_K^2}{T}.
\]
This is sharper than the displayed target and therefore implies
\(O_p(\zeta_K^2K/T)\), because \(K\geq1\).  All inequalities hold on a single
uniform Gram event whose probability tends to one.
\end{proof}


\section{Known-Partition Primitive Wilks}

Set
\[
    \bm v_0=\bm M_{K,\Lambda_0}\bm w,
    \qquad
    A_T^{(q)}=T^{-1/2}\sum_{t=1}^TU_tv_{t,0},
    \qquad
    V_T^{(q)}=T^{-1}\sum_{t=1}^TU_t^2v_{t,0}^2.
\]

\begin{lemma}[Oracle martingale CLT target]
\label{lem:wb-oracle-clt}
Under Assumptions \ref{ass:prim-stoch}, \ref{ass:prim-score-admissible},
\ref{ass:prim-bounded}, \ref{ass:prim-gram}, and
\ref{ass:sol-oracle-envelope}, suppose
\[
    T^{-1}\sum_{t=1}^T \sigma_t^2v_{t,0}^2\xrightarrow{p}\eta_q^2,
    \qquad
    \Pp(\eta_q\geq \sqrt{c_\eta})=1
\]
for some \(c_\eta>0\). Then
\[
    A_T^{(q)}\xrightarrow{\mathrm{st}}\eta_qZ,
\]
where \(Z\sim N(0,1)\) is independent of the limiting sigma-field generating \(\eta_q\).
\end{lemma}


\begin{proof}[Detailed solution]
Let \(\delta_0=\min\{\delta,2\}\).  First use the predictable route.  Put
\[
    X_{t,T}=T^{-1/2}U_tv_{t,0}^\ast.
\]
Then \(X_{t,T}\) is a martingale-difference array.  By Assumption
\ref{ass:sol-oracle-envelope} and the variance-replacement condition,
\begin{align*}
    \sum_{t=1}^T\E(X_{t,T}^2\mid\mathcal F_{t-1})
    &=T^{-1}\sum_{t=1}^T\sigma_t^2(v_{t,0}^\ast)^2\\
    &=T^{-1}\sum_{t=1}^T\sigma_t^2v_{t,0}^2+o_p(1)
      \xrightarrow{p}\eta_q^2.
\end{align*}
For the conditional Lyapunov term, the conditional moment bound gives
\begin{align*}
    \sum_{t=1}^T
      \E(|X_{t,T}|^{2+\delta_0}\mid\mathcal F_{t-1})
    &\leq
      C T^{-(1+\delta_0/2)}
      \sum_{t=1}^T|v_{t,0}^\ast|^{2+\delta_0}\\
    &\leq
      C b_K^{2+\delta_0}T^{-\delta_0/2}
      =o_p(1).
\end{align*}
The martingale triangular-array stable CLT therefore yields
\[
    T^{-1/2}\sum_{t=1}^TU_tv_{t,0}^\ast
    \xrightarrow{\mathrm{st}}\eta_qZ,
\]
where \(Z\) is standard normal and independent of the limiting sigma-field
that contains \(\eta_q\).  Assumption \ref{ass:prim-score-admissible} gives
\[
    T^{-1/2}\sum_{t=1}^TU_t(v_{t,0}-v_{t,0}^\ast)=o_p(1).
\]
Stable Slutsky therefore gives the asserted limit for \(A_T^{(q)}\).  Under
Route A the same argument is exactly the admissible martingale CLT postulated
there, so no replacement is required.
\end{proof}


\begin{lemma}[Oracle denominator LLN target]
\label{lem:wb-oracle-lln}
Under the primitive stochastic and score-admissibility assumptions and
Assumption \ref{ass:sol-oracle-envelope},
\[
    V_T^{(q)}-\eta_q^2=o_p(1).
\]
\end{lemma}


\begin{proof}[Detailed solution]
Again use the predictable approximation; the direct predictable case is the
special case \(v_{t,0}^\ast=v_{t,0}\).  Decompose
\begin{align*}
    T^{-1}\sum_{t=1}^TU_t^2(v_{t,0}^\ast)^2
    &=T^{-1}\sum_{t=1}^T\sigma_t^2(v_{t,0}^\ast)^2\\
    &\quad+T^{-1}\sum_{t=1}^T
       (U_t^2-\sigma_t^2)(v_{t,0}^\ast)^2.
\end{align*}
The first term converges in probability to \(\eta_q^2\) by Assumption
\ref{ass:sol-oracle-envelope} and the variance convergence in Lemma
\ref{lem:wb-oracle-clt}.

For the second term, set
\[
    D_{t,T}=(U_t^2-\sigma_t^2)(v_{t,0}^\ast)^2,
    \qquad p=1+\delta_0/2\in(1,2].
\]
Then \(D_{t,T}\) is a martingale difference and, using
\(2p=2+\delta_0\),
\[
    \E(|D_{t,T}|^p\mid\mathcal F_{t-1})
    \leq C|v_{t,0}^\ast|^{2p}.
\]
The martingale von Bahr--Esseen inequality gives
\begin{align*}
    \E\left|
      T^{-1}\sum_{t=1}^TD_{t,T}
    \right|^p
    &\leq C T^{-p}\sum_{t=1}^T\E|D_{t,T}|^p\\
    &\leq C b_K^{2+\delta_0}T^{-\delta_0/2}=o(1),
\end{align*}
where the usual localization to the event
\(\max_t|v_{t,0}^\ast|\leq Cb_K\) is understood.  Hence the martingale term is
\(o_p(1)\).  Finally, Assumption \ref{ass:prim-score-admissible} supplies
\[
    T^{-1}\sum_{t=1}^TU_t^2
       \{v_{t,0}^2-(v_{t,0}^\ast)^2\}=o_p(1),
\]
so replacing \(v_{t,0}^\ast\) by \(v_{t,0}\) proves
\(V_T^{(q)}=\eta_q^2+o_p(1)\).
\end{proof}


\begin{lemma}[Oracle maximum bound target]
\label{lem:wb-oracle-max}
Under the primitive moment and leverage conditions and Assumption
\ref{ass:sol-oracle-envelope},
\[
    \max_{t\leq T}|U_tv_{t,0}|=o_p(\sqrt T).
\]
\end{lemma}


\begin{proof}[Detailed solution]
Lemma \ref{lem:wb-leverage} gives
\[
    \max_{t\leq T}|v_{t,0}|=O_p(b_K).
\]
Fix \(C<\infty\) and work on the event
\(\max_t|v_{t,0}|\leq Cb_K\).  For every \(\varepsilon>0\), the union bound
and the \((2+\delta_0)\)-moment bound imply
\begin{align*}
    \Pp\left(
      \max_{t\leq T}|U_tv_{t,0}|>\varepsilon\sqrt T
    \right)
    &\leq
      \sum_{t=1}^T
      \Pp\left(
        |U_t|>\frac{\varepsilon\sqrt T}{Cb_K}
      \right)+o(1)\\
    &\leq
      C_\varepsilon T
      \left(\frac{b_K}{\sqrt T}\right)^{2+\delta_0}+o(1)\\
    &=C_\varepsilon b_K^{2+\delta_0}T^{-\delta_0/2}+o(1)
      \longrightarrow0.
\end{align*}
Thus \(\max_t|U_tv_{t,0}|/\sqrt T\to_p0\).  Equivalently, a sufficient rate
is
\[
    b_KT^{1/(2+\delta_0)-1/2}\to0,
\]
which is exactly \eqref{eq:sol-envelope-rate} written after taking a
\((2+\delta_0)\)-th root.
\end{proof}


\begin{lemma}[Feasible true-partition replacement]
\label{lem:wb-feasible-replacement}
Under Assumptions \ref{ass:prim-stoch}, \ref{ass:prim-sieve},
\ref{ass:prim-bounded}, \ref{ass:prim-gram}, and
\ref{ass:sol-oracle-envelope},
\begin{align}
    T^{-1/2}\sum_{t=1}^T\widehat Z_t(\beta_0,\Lambda_0)
    &=
    T^{-1/2}\sum_{t=1}^TU_tv_{t,0}+o_p(1),
    \\
    T^{-1}\sum_{t=1}^T\widehat Z_t(\beta_0,\Lambda_0)^2
    &=
    T^{-1}\sum_{t=1}^TU_t^2v_{t,0}^2+o_p(1),
    \\
    \max_{t\leq T}|\widehat Z_t(\beta_0,\Lambda_0)|
    &=o_p(\sqrt T).
\end{align}
\end{lemma}


\begin{proof}[Detailed solution]
Let \(\bm Q_0=\bm Q_{K,\Lambda_0}\),
\(\bm\Pi_0=\bm I_T-\bm M_{K,\Lambda_0}\), and
\(\bm v_0=\bm M_{K,\Lambda_0}\bm w\).  The sieve representation gives
\[
    \bm Y-\beta_0\bm X_{lag}
    =\bm Q_0\bm c_0+\bm r_0+\bm U.
\]
By Lemma \ref{lem:wb-projection},
\[
    \widehat{\bm u}(\beta_0,\Lambda_0)
    =\bm M_{K,\Lambda_0}(\bm r_0+\bm U).
\]
Define
\[
    \bm e_0=\bm M_{K,\Lambda_0}\bm r_0,
    \qquad
    \bm h_0=\bm\Pi_0\bm U.
\]
Then, componentwise,
\begin{equation}
    \widehat Z_t(\beta_0,\Lambda_0)
    =U_tv_{t,0}+(e_{t,0}-h_{t,0})v_{t,0}.
    \label{eq:sol-feasible-score-decomp}
\end{equation}

We first establish the projection-noise bounds rather than assuming them.
Because each row of \(\bm Q_0\) is predictable, martingale orthogonality and
the Gram upper bound imply
\[
    \|\bm Q_0'\bm U\|=O_p(\sqrt{Td_K}).
\]
Indeed, after localization to the Gram event,
\[
    \E\|\bm Q_0'\bm U\|^2
    =\sum_{t=1}^T\E\{\sigma_t^2\|\bm q_{t,0}\|^2\}
    \leq C\E\operatorname{tr}(\bm Q_0'\bm Q_0)
    =O(Td_K).
\]
Since \(\|(\bm Q_0'\bm Q_0)^\dagger\|=O_p(T^{-1})\),
\begin{align*}
    \|\bm h_0\|_T^2
    &=T^{-1}\bm U'\bm\Pi_0\bm U\\
    &\leq
      \frac{C}{T^2}\|\bm Q_0'\bm U\|^2
      =O_p(d_K/T),
\end{align*}
and, row by row,
\[
    \|\bm h_0\|_\infty
    \leq
    \zeta_K\|(\bm Q_0'\bm Q_0)^\dagger\|
             \|\bm Q_0'\bm U\|
    =O_p\left(\zeta_K\sqrt{d_K/T}\right).
\]

For the approximation term, orthogonal projection is a contraction, so
\[
    \|\bm e_0\|_T\leq\|\bm r_0\|_T=O_p(a_K).
\]
Furthermore,
\[
    \|(\bm Q_0'\bm Q_0)^\dagger\bm Q_0'\bm r_0\|
    \leq \frac{C}{T}\|\bm Q_0\|\,\|\bm r_0\|
    =O_p(a_K),
\]
which gives
\[
    \|\bm e_0\|_\infty
    \leq \|\bm r_0\|_\infty+\|\bm\Pi_0\bm r_0\|_\infty
    =O_p\{a_K(1+\zeta_K)\}.
\]
Combining these bounds with
\(\|\bm v_0\|_\infty=O_p(b_K)\),
\begin{align*}
    \|(\bm e_0-\bm h_0)\odot\bm v_0\|_T
    &\leq O_p\left(a_Kb_K+b_K\sqrt{d_K/T}\right)=o_p(1),\\
    \|(\bm e_0-\bm h_0)\odot\bm v_0\|_\infty
    &\leq O_p\left(a_Kb_K^2+\zeta_Kb_K\sqrt{d_K/T}\right)
      =o_p(\sqrt T).
\end{align*}
The first conclusion uses Assumption \ref{ass:prim-sieve}; the second uses
both Assumptions \ref{ass:prim-sieve} and
\ref{ass:sol-oracle-envelope}.

For the numerator, symmetry and idempotence give an exact cancellation:
\begin{align*}
    T^{-1/2}\sum_t\widehat Z_t(\beta_0,\Lambda_0)
    &=T^{-1/2}(\bm U+\bm r_0)'
      \bm M_{K,\Lambda_0}^2\bm w\\
    &=T^{-1/2}\bm U'\bm v_0
      +T^{-1/2}\bm r_0'\bm v_0\\
    &=T^{-1/2}\bm U'\bm v_0+o_p(1).
\end{align*}
For the denominator, let
\(\bm d_0=(\bm e_0-\bm h_0)\odot\bm v_0\).  Then
\begin{align*}
    \left|
      T^{-1}\sum_t\widehat Z_t^2
      -T^{-1}\sum_tU_t^2v_{t,0}^2
    \right|
    &\leq
      2\|\bm U\odot\bm v_0\|_T\|\bm d_0\|_T
      +\|\bm d_0\|_T^2\\
    &=o_p(1),
\end{align*}
because Lemma \ref{lem:wb-oracle-lln} makes the first norm \(O_p(1)\).
Finally, \eqref{eq:sol-feasible-score-decomp}, Lemma
\ref{lem:wb-oracle-max}, and \(\|\bm d_0\|_\infty=o_p(\sqrt T)\) prove the
maximum assertion.
\end{proof}


\begin{theorem}[Primitive known-partition Wilks target]
\label{thm:wb-known-wilks}
Under the primitive assumptions in this section, Assumption
\ref{ass:sol-oracle-envelope}, the stated variance convergence, and the
convex-hull condition for \(\widehat Z_t(\beta_0,\Lambda_0)\),
\[
    \ell_T(\beta_0,\Lambda_0)\Rightarrow\chi_1^2.
\]
\end{theorem}


\begin{proof}[Detailed solution]
Set
\[
    A_T=T^{-1/2}\sum_{t=1}^T
       \widehat Z_t(\beta_0,\Lambda_0),
    \qquad
    B_T=T^{-1}\sum_{t=1}^T
       \widehat Z_t(\beta_0,\Lambda_0)^2.
\]
Lemma \ref{lem:wb-feasible-replacement}, followed by Lemmas
\ref{lem:wb-oracle-clt} and \ref{lem:wb-oracle-lln}, gives the joint limits
\[
    A_T\xrightarrow{\mathrm{st}}\eta_qZ,
    \qquad
    B_T\xrightarrow{p}\eta_q^2.
\]
Because \(\eta_q^2\geq c_\eta\) almost surely, \(B_T\) is bounded away from
zero with probability tending to one.  Lemmas \ref{lem:wb-oracle-max} and
\ref{lem:wb-feasible-replacement} give
\[
    \max_t|\widehat Z_t(\beta_0,\Lambda_0)|=o_p(\sqrt T).
\]
The convex-hull assumption permits application of Lemma
\ref{lem:sol-scalar-el}; hence
\[
    \ell_T(\beta_0,\Lambda_0)
    =\frac{A_T^2}{B_T}+o_p(1).
\]
Stable Slutsky and positivity of \(\eta_q\) imply
\[
    \frac{A_T^2}{B_T}
    \Rightarrow
    \frac{\eta_q^2Z^2}{\eta_q^2}=Z^2\sim\chi_1^2.
\]
This proves the theorem.
\end{proof}


\section{One-Break Estimated-Partition Bridge}
\label{sec:one-break-bridge}

In this section only, take \(q_0=1\), \(\Lambda_0=(k_0)\), and \(\widetilde\Lambda=(\widetilde k)\). Let
\[
    \Delta_M=\bm M_{K,\widetilde\Lambda}-\bm M_{K,\Lambda_0},
    \qquad
    \mathcal M_T=\{t:j_{\widetilde\Lambda}(t)\neq j_{\Lambda_0}(t)\}.
\]

\begin{lemma}[Boundary-set size]
\label{lem:wb-boundary-size}
If \(|\widetilde k-k_0|=O_p(r_T)\), then
\[
    |\mathcal M_T|=O_p(r_T).
\]
\end{lemma}


\begin{proof}[Detailed solution]
Suppose first that \(\widetilde k\geq k_0\).  For
\(t\leq k_0\), both partitions assign observation \(t\) to the left regime;
for \(t>\widetilde k\), both assign it to the right regime.  The assignments
differ exactly for
\[
    k_0<t\leq\widetilde k.
\]
Thus \(|\mathcal M_T|=\widetilde k-k_0\).  If
\(\widetilde k<k_0\), the differing observations are exactly
\(\widetilde k<t\leq k_0\), and again
\(|\mathcal M_T|=|\widetilde k-k_0|\).  Therefore
\[
    |\mathcal M_T|=|\widetilde k-k_0|=O_p(r_T).
\]
\end{proof}


\begin{lemma}[Nuisance part of partition perturbation]
\label{lem:wb-bridge-mu}
Under bounded nuisance functions, bounded weights, uniform Gram stability,
Assumption \ref{ass:sol-local-bias}, and
\(|\widetilde k-k_0|=O_p(r_T)\),
\[
    T^{-1/2}\bm\mu'\Delta_M\bm w
    =
    O_p\!\left(\frac{(1+\zeta_K)r_T}{\sqrt T}\right)+o_p(1).
\]
In particular, this term is \(o_p(1)\) if \((1+\zeta_K)r_T/\sqrt T\to0\).
\end{lemma}


\begin{proof}[Detailed solution]
The original primitive assumptions alone do not control the term generated
by moving the sieve approximation error from one projection to another;
Assumption \ref{ass:sol-local-bias} supplies exactly that missing control.
We now prove the stated rate.

Work first on the event \(\widetilde k=k_0+h\) with \(h\geq0\), and let
\(H=\{k_0+1,\ldots,k_0+h\}\).  Let \(\bm c_0\) and \(\bm c_1\) be the two
true-regime sieve coefficients.  Form a coefficient vector for the candidate
partition by assigning \(\bm c_0\) to its left block and \(\bm c_1\) to its
right block.  Then
\[
    \bm\mu=\bm Q_{K,\widetilde\Lambda}\bm c^{(\widetilde k)}
            +\bm d_H+\bm r_0,
\]
where \(\bm d_H\) is supported on \(H\) and
\[
    d_{H,t}=\bm P_K(W_{t-1})'(\bm c_1-\bm c_0),
    \qquad t\in H.
\]
Because, pointwise for \(t\in H\),
\[
    \bm P_K(W_{t-1})'(\bm c_1-\bm c_0)
    =m_1(W_{t-1})-m_0(W_{t-1})
      -\{r_{1,K,t}-r_{0,K,t}\},
\]
Assumptions \ref{ass:prim-bounded} and \ref{ass:prim-sieve} imply
\(\|\bm d_H\|_\infty=O_p(1)\).  Orthogonality of the candidate block sieve
therefore gives
\[
    \bm\mu'\bm M_{K,\widetilde\Lambda}\bm w
    =\bm d_H'\bm M_{K,\widetilde\Lambda}\bm w
      +\bm r_0'\bm M_{K,\widetilde\Lambda}\bm w.
\]
At the true partition,
\[
    \bm\mu'\bm M_{K,\Lambda_0}\bm w
    =\bm r_0'\bm M_{K,\Lambda_0}\bm w.
\]
Subtracting and using Lemma \ref{lem:wb-leverage},
\begin{align*}
    T^{-1/2}|\bm\mu'\Delta_M\bm w|
    &\leq
      T^{-1/2}\|\bm d_H\|_1
      \|\bm M_{K,\widetilde\Lambda}\bm w\|_\infty\\
    &\quad+
      T^{-1/2}|\bm r_0'\Delta_M\bm w|\\
    &=O_p\left(\frac{b_Kh}{\sqrt T}\right)+o_p(1).
\end{align*}
The case \(\widetilde k<k_0\) is identical after exchanging the two regimes.
Since \(h=O_p(r_T)\), the displayed target follows.

For completeness, here is the promised sufficient verification of
Assumption \ref{ass:sol-local-bias}.  Let
\(\bm v_k=\bm M_{K,k}\bm w\).  The coefficient-update calculation in the
proof of Lemma \ref{lem:wb-bridge-u} shows that, for
\(h=|k-k_0|\leq Cr_T\),
\[
    \|\bm v_k-\bm v_{k_0}\|_1
    =O_p\{h b_K(1+\zeta_K^2)\}.
\]
Consequently,
\[
    T^{-1/2}|\bm r_0'(\bm v_k-\bm v_{k_0})|
    \leq
    T^{-1/2}\|\bm r_0\|_\infty
               \|\bm v_k-\bm v_{k_0}\|_1
    =O_p\left(
       \frac{a_Kb_K(1+\zeta_K^2)h}{\sqrt T}
      \right),
\]
which is uniformly \(o_p(1)\) under
\eqref{eq:sol-local-bias-sufficient}.
\end{proof}


\begin{lemma}[Stochastic part of partition perturbation]
\label{lem:wb-bridge-u}
Under the martingale error, moment, bounded-basis, and uniform Gram
assumptions, Assumption \ref{ass:sol-local-series}, and
\(b_Kr_T/\sqrt T\to0\),
\[
    T^{-1/2}\bm U'\Delta_M\bm w
    =
    O_p\!\left(\zeta_K\sqrt{\frac{K r_T}{T}}\right)+o_p(1).
\]
In particular, this term is \(o_p(1)\) if \(\zeta_K\sqrt{K r_T/T}\to0\).
\end{lemma}


\begin{proof}[Detailed solution]
It is enough to work uniformly on
\(|k-k_0|\leq Cr_T\), because the localization event has probability tending
to one for sufficiently large fixed \(C\).  Consider \(k=k_0+h>k_0\) and
write
\[
    A=\{1,\ldots,k_0\},\quad
    H=\{k_0+1,\ldots,k_0+h\},\quad
    B=\{k_0+h+1,\ldots,T\},\quad R=H\cup B.
\]
For a segment \(J\), put
\[
    \bm G_J=\bm P_J'\bm P_J,
    \qquad
    \widehat{\bm\gamma}_J=\bm G_J^{-1}\bm P_J'\bm w_J.
\]
Assumption \ref{ass:sol-local-series} and bounded \(\bm w\) imply
\(\|\widehat{\bm\gamma}_J\|=O_p(1)\) for every long segment.  The exact update
identity
\[
    \widehat{\bm\gamma}_{A\cup H}-\widehat{\bm\gamma}_A
    =\bm G_{A\cup H}^{-1}
       \bm P_H'(\bm w_H-\bm P_H\widehat{\bm\gamma}_A)
\]
gives
\[
    \|\widehat{\bm\gamma}_{A\cup H}-
       \widehat{\bm\gamma}_A\|
    =O_p\left(\frac{h\zeta_Kb_K}{T}\right).
\]
Removing \(H\) from the true right segment gives, by the same algebra,
\[
    \|\widehat{\bm\gamma}_{B}-
       \widehat{\bm\gamma}_{R}\|
    =O_p\left(\frac{h\zeta_Kb_K}{T}\right).
\]

The true and candidate residualized weights are segmentwise
\[
    \bm v_0=(\bm w_A-\bm P_A\widehat{\bm\gamma}_A,
             \bm w_R-\bm P_R\widehat{\bm\gamma}_R),
\]
\[
    \bm v_k=(\bm w_{A\cup H}-
              \bm P_{A\cup H}\widehat{\bm\gamma}_{A\cup H},
             \bm w_B-\bm P_B\widehat{\bm\gamma}_B).
\]
After cancelling the \(\bm w_H\) terms, one obtains the exact decomposition
\begin{align}
    \bm U'(\bm v_k-\bm v_0)
    &= (\bm P_A'\bm U_A)'
       (\widehat{\bm\gamma}_A-
        \widehat{\bm\gamma}_{A\cup H})
       \nonumber\\
    &\quad +(\bm P_B'\bm U_B)'
       (\widehat{\bm\gamma}_R-
        \widehat{\bm\gamma}_{B})
       \nonumber\\
    &\quad +(\bm P_H'\bm U_H)'
       (\widehat{\bm\gamma}_R-
        \widehat{\bm\gamma}_{A\cup H}).
    \label{eq:sol-u-bridge-decomp}
\end{align}
Because \(\bm P_K(W_{t-1})U_t\) is a vector martingale difference, martingale
isometry and the Gram bounds yield
\[
    \|\bm P_A'\bm U_A\|=O_p(\sqrt{TK}),
    \qquad
    \sup_{h\leq Cr_T}\|\bm P_H'\bm U_H\|
      =O_p(\zeta_K\sqrt{r_T}).
\]
The same bound for \(\bm P_B'\bm U_B\) follows by writing it as the fixed
right-segment sum minus the local \(H\)-sum.  Insert these bounds into
\eqref{eq:sol-u-bridge-decomp}.  Uniformly for \(h\leq Cr_T\),
\begin{align*}
    T^{-1/2}|\bm U'(\bm v_k-\bm v_0)|
    &=O_p\left(\zeta_K\sqrt{\frac{h}{T}}\right)
      +O_p\left(\frac{h\zeta_Kb_K\sqrt K}{T}\right)
      +o_p(1).
\end{align*}
The second displayed order divided by
\(\zeta_K\sqrt{Kh/T}\) equals \(b_K\sqrt{h/T}\).  Under the bridge rate
\(b_Kr_T/\sqrt T\to0\), this ratio is \(o(1)\), uniformly for
\(h\leq Cr_T\).  Since \(K\geq1\), the first order is bounded by
\(\zeta_K\sqrt{Kh/T}\).  Thus
\[
    T^{-1/2}\bm U'\Delta_M\bm w
    =O_p\left(\zeta_K\sqrt{\frac{Kr_T}{T}}\right)+o_p(1).
\]
The case \(k<k_0\) is obtained by reversing left and right.  Notice that the
proof is uniform over deterministic local candidates, so it remains valid at
the data-selected \(\widetilde k\).
\end{proof}


\begin{lemma}[Denominator perturbation]
\label{lem:wb-bridge-v}
Under the conditions of Lemmas \ref{lem:wb-bridge-mu} and
\ref{lem:wb-bridge-u}, and under \(\zeta_K^2Kr_T/T\to0\),
\[
    V_T(\widetilde\Lambda)-V_T(\Lambda_0)
    =
    O_p\!\left(
      \frac{(1+\zeta_K)^2r_T}{T}
      +
      \frac{\zeta_K^2Kr_T}{T}
    \right)+o_p(1)
    =o_p(1).
\]
\end{lemma}


\begin{proof}[Detailed solution]
Put \(\bm v_k=\bm M_{K,k}\bm w\).  We first reduce the feasible denominator
to the oracle array \(U_t v_{t,k}\), uniformly for
\(|k-k_0|\leq Cr_T\).  As in the proof of Lemma \ref{lem:wb-bridge-mu},
\[
    \bm\mu=\bm Q_{K,k}\bm c^{(k)}+\bm d_H+\bm r_0,
\]
where \(\bm d_H\) is supported on the boundary set and is uniformly bounded.
Writing \(\bm\Pi_k=\bm I_T-\bm M_{K,k}\),
\[
    \widehat{\bm u}(\beta_0,k)
    =\bm U-\bm\Pi_k\bm U+
      \bm M_{K,k}\bm r_0+\bm M_{K,k}\bm d_H.
\]
Hence
\[
    \widehat{\bm Z}(\beta_0,k)
    =\bm U\odot\bm v_k+\bm\rho_k,
\]
where
\[
    \|\bm\rho_k\|_T
    \leq \|\bm v_k\|_\infty
       \left(
        \|\bm\Pi_k\bm U\|_T+
        \|\bm M_{K,k}\bm r_0\|_T+
        \|\bm M_{K,k}\bm d_H\|_T
       \right).
\]
The projection-noise proof used in Lemma
\ref{lem:wb-feasible-replacement}, now uniformly over local candidates, gives
\(\|\bm\Pi_k\bm U\|_T=O_p(\sqrt{K/T})\).  Projection contraction gives
\[
    \|\bm M_{K,k}\bm r_0\|_T=O_p(a_K),
    \qquad
    \|\bm M_{K,k}\bm d_H\|_T=O_p(\sqrt{h/T}).
\]
Together with \(\|\bm v_k\|_\infty=O_p(b_K)\), the sieve and bridge rates
imply
\begin{equation}
    \sup_{|k-k_0|\leq Cr_T}\|\bm\rho_k\|_T=o_p(1).
    \label{eq:sol-rho-local}
\end{equation}
Therefore it remains to compare
\[
    V_{U,T}(k)=T^{-1}\sum_{t=1}^TU_t^2v_{t,k}^2.
\]

Let \(h=|k-k_0|\).  The coefficient-update bounds proved in Lemma
\ref{lem:wb-bridge-u} imply
\[
    |v_{t,k}-v_{t,0}|
    \leq
    \begin{cases}
      Cb_K, & t\in\mathcal M_T,\\[2mm]
      C h\zeta_K^2b_K/T, & t\notin\mathcal M_T.
    \end{cases}
\]
Also \(\max_t(|v_{t,k}|+|v_{t,0}|)=O_p(b_K)\).  Since \(T^{-1}\sum_tU_t^2=O_p(1)\), it remains only to justify the
boundary quadratic bound.  Put
\(p=1+\delta_0/2>1\).  Applied to the centered sequence
\(U_t^2-\E(U_t^2\mid\mathcal F_{t-1})\), the martingale maximal inequality on
dyadic blocks gives
\[
    \sup_{1\leq h\leq Cr_T}
    \frac{1}{h}\left|
      \sum_{t\in H_h}
       \{U_t^2-\E(U_t^2\mid\mathcal F_{t-1})\}
    \right|=O_p(1),
\]
because the block-tail probabilities are bounded by a summable geometric
series of order \(2^{-j(p-1)}\).  The conditional variances are bounded, so
uniformly for a boundary interval \(H_h\) of length \(h\leq Cr_T\),
\[
    \sum_{t\in H_h}U_t^2=O_p(h).
\]
Consequently,
\begin{align*}
    |V_{U,T}(k)-V_{U,T}(k_0)|
    &\leq T^{-1}\sum_tU_t^2
       |v_{t,k}-v_{t,0}|(|v_{t,k}|+|v_{t,0}|)\\
    &=O_p\left(\frac{b_K^2h}{T}\right)
      +O_p\left(\frac{h\zeta_K^2b_K^2}{T}\right)
      +O_p\left(\frac{h^2\zeta_K^4b_K^2}{T^2}\right).
\end{align*}
Assumption \ref{ass:sol-local-series} gives \(b_K^2=O(K)\).  Moreover,
\(h\zeta_K^2/T=o(1)\) follows from
\(\zeta_K^2Kr_T/T\to0\).  Consequently,
\[
    V_{U,T}(k)-V_{U,T}(k_0)
    =O_p\left(
       \frac{b_K^2h}{T}+
       \frac{\zeta_K^2Kh}{T}
      \right).
\]
This bound is \(o_p(1)\), so
\(\|\bm U\odot\bm v_k\|_T=O_p(1)\) follows from the true-partition oracle
LLN.  Now \eqref{eq:sol-rho-local} and the square-difference inequality give
\[
    V_T(k)-V_{U,T}(k)=o_p(1)
\]
uniformly locally, and the same holds at \(k_0\).  Taking
\(h=O_p(r_T)\) proves
\[
    V_T(\widetilde\Lambda)-V_T(\Lambda_0)
    =O_p\left(
      \frac{(1+\zeta_K)^2r_T}{T}
      +\frac{\zeta_K^2Kr_T}{T}
    \right)+o_p(1).
\]
\end{proof}


\begin{theorem}[One-break estimated-partition score equivalence]
\label{thm:wb-one-break-equivalence}
Suppose \(q_0=1\), Assumption \ref{ass:prim-localization-target} holds with \(r_T=o(\sqrt T)\), and
\[
    \zeta_K\sqrt{\frac{Kr_T}{T}}\to0,
    \qquad
    \frac{(1+\zeta_K)r_T}{\sqrt T}\to0,
    \qquad
    \frac{\zeta_K^2Kr_T}{T}\to0.
\]
If Lemmas \ref{lem:wb-bridge-mu}, \ref{lem:wb-bridge-u}, and \ref{lem:wb-bridge-v} hold, then
\[
    S_T(\widetilde\Lambda)-S_T(\Lambda_0)=o_p(1),
    \qquad
    V_T(\widetilde\Lambda)-V_T(\Lambda_0)=o_p(1).
\]
\end{theorem}


\begin{proof}[Detailed solution]
For every partition \(\Lambda\), symmetry and idempotence imply
\[
    S_T(\Lambda)
    =T^{-1/2}(\bm U+\bm\mu)'
       \bm M_{K,\Lambda}\bm w.
\]
Therefore
\[
    S_T(\widetilde\Lambda)-S_T(\Lambda_0)
    =T^{-1/2}\bm\mu'\Delta_M\bm w
      +T^{-1/2}\bm U'\Delta_M\bm w.
\]
Lemma \ref{lem:wb-bridge-mu} bounds the first term by
\[
    O_p\left(\frac{b_Kr_T}{\sqrt T}\right)+o_p(1)=o_p(1),
\]
and Lemma \ref{lem:wb-bridge-u} bounds the second by
\[
    O_p\left(\zeta_K\sqrt{\frac{Kr_T}{T}}\right)+o_p(1)=o_p(1).
\]
This proves the numerator equivalence.

For the denominator, Lemma \ref{lem:wb-bridge-v} gives
\[
    V_T(\widetilde\Lambda)-V_T(\Lambda_0)
    =O_p\left(
       \frac{b_K^2r_T}{T}+
       \frac{\zeta_K^2Kr_T}{T}
      \right)+o_p(1).
\]
The second displayed bridge rate makes the second term \(o_p(1)\).  For the
first term, if \(r_T\geq1\),
\[
    \frac{b_K^2r_T}{T}
    =\frac{1}{r_T}
      \left(\frac{b_Kr_T}{\sqrt T}\right)^2=o(1).
\]
Hence the denominator difference is also \(o_p(1)\).
\end{proof}


\begin{theorem}[One-break primitive estimated-partition Wilks target]
\label{thm:wb-one-break-wilks}
Under the assumptions of Theorem \ref{thm:wb-known-wilks}, the convex-hull condition for the estimated-partition score array, and Theorem \ref{thm:wb-one-break-equivalence},
\[
    \ell_T(\beta_0,\widetilde\Lambda)\Rightarrow\chi_1^2.
\]
\end{theorem}


\begin{proof}[Detailed solution]
Let \(A_T(\Lambda)=S_T(\Lambda)\) and \(B_T(\Lambda)=V_T(\Lambda)\).
Theorem \ref{thm:wb-one-break-equivalence} gives
\[
    A_T(\widetilde\Lambda)=A_T(\Lambda_0)+o_p(1),
    \qquad
    B_T(\widetilde\Lambda)=B_T(\Lambda_0)+o_p(1).
\]
The known-partition proof gives
\[
    A_T(\Lambda_0)\xrightarrow{\mathrm{st}}\eta_qZ,
    \qquad
    B_T(\Lambda_0)\xrightarrow{p}\eta_q^2.
\]
We must also verify the maximum condition; numerator and denominator
 equivalence alone would not imply it.  The local representation in the proof
of Lemma \ref{lem:wb-bridge-v} gives
\[
    \widehat{\bm Z}(\beta_0,\widetilde\Lambda)
    =\bm U\odot\bm v_{\widetilde k}+\bm\rho_{\widetilde k}.
\]
The envelope rate implies
\(\max_t|U_tv_{t,\widetilde k}|=o_p(\sqrt T)\).  For the remainder, the
uniform projection calculation gives
\[
    \|\bm\Pi_{\widetilde k}\bm U\|_\infty
    =O_p\left(\zeta_K\sqrt{K/T}\right),
\]
while the sieve calculation gives
\[
    \|\bm M_{K,\widetilde k}\bm r_0\|_\infty
    =O_p\{a_K(1+\zeta_K)\}.
\]
Finally, because \(\bm d_H\) has bounded entries and support of size
\(O_p(r_T)\),
\[
    \|\bm M_{K,\widetilde k}\bm d_H\|_\infty
    =O_p\left(1+\frac{r_T\zeta_K^2}{T}\right)=O_p(1).
\]
Multiplication by \(\|\bm v_{\widetilde k}\|_\infty=O_p(b_K)\), followed by
the sieve, bridge, and envelope rates, gives
\(\|\bm\rho_{\widetilde k}\|_\infty=o_p(\sqrt T)\).  Hence
\[
    \max_t|\widehat Z_t(\beta_0,\widetilde\Lambda)|=o_p(\sqrt T).
\]
The estimated-partition convex-hull condition and Lemma
\ref{lem:sol-scalar-el} now yield
\[
    \ell_T(\beta_0,\widetilde\Lambda)
    =\frac{A_T(\widetilde\Lambda)^2}
           {B_T(\widetilde\Lambda)}+o_p(1).
\]
Substitution of the two equivalences reduces this ratio to the
true-partition ratio plus \(o_p(1)\), and Theorem
\ref{thm:wb-known-wilks} gives the \(\chi_1^2\) limit.
\end{proof}


\section{Break-Date Localization}
\label{sec:break-localization}

\begin{lemma}[Deterministic RSS drift target]
\label{lem:wb-rss-drift}
Under Assumption \ref{ass:sol-localization}, there exists
\(c_\Delta>0\) such that, uniformly over minimum-spaced one-break candidates
\(k\),
\[
    \E\{RSS_K(\beta_0,k)-RSS_K(\beta_0,k_0)\mid \mathcal D_T\}
    \geq
    c_\Delta |k-k_0|\Delta_T^2
    -
    o_p(|k-k_0|\Delta_T^2),
\]
where \(\mathcal D_T\) denotes the design sigma-field generated by \(X,W\) and the sieve basis.
\end{lemma}


\begin{proof}[Detailed solution]
The conclusion is not valid from the original martingale assumptions alone;
we use Assumption \ref{ass:sol-localization}.  Let
\[
    \bm y_0=\bm Y-\beta_0\bm X_{lag}=\bm\mu+\bm U,
    \qquad
    \bm A_k=\bm M_{K,k}-\bm M_{K,k_0}.
\]
Because the profile least-squares cost is the residual quadratic form,
\[
    RSS_K(\beta_0,k)=\bm y_0'\bm M_{K,k}\bm y_0.
\]
Therefore
\begin{align*}
    \Delta RSS_K(k)
    &=\bm\mu'\bm A_k\bm\mu
      +2\bm\mu'\bm A_k\bm U
      +\bm U'\bm A_k\bm U.
\end{align*}
Conditional centering and conditional homoskedasticity give
\[
    \E\{\Delta RSS_K(k)\mid\mathcal D_T\}
    =\bm\mu'\bm A_k\bm\mu
      +\sigma^2\operatorname{tr}(\bm A_k).
\]
Every candidate and true one-break sieve matrix has rank \(2K\), so
\[
    \operatorname{tr}(\bm A_k)
    =\operatorname{tr}(\bm M_{K,k})-
      \operatorname{tr}(\bm M_{K,k_0})=0.
\]
The profile-identification condition
\eqref{eq:sol-profile-identification} now gives
\[
    \E\{\Delta RSS_K(k)\mid\mathcal D_T\}
    \geq c_\Delta|k-k_0|\Delta_T^2
       -o_p\{|k-k_0|\Delta_T^2\}.
\]

To see concretely where the linear drift comes from, consider the exact-sieve
case and \(k=k_0+h>k_0\).  Use the sets \(A,H,B\) from the bridge proof and
write \(\bm\delta=\bm c_1-\bm c_0\).  The candidate right segment fits
exactly.  On the candidate left segment, after writing the common coefficient
as \(\bm c_0+\bm a\), the deterministic loss is
\[
    L_h=\min_{\bm a}
       \{\bm a'\bm G_A\bm a+
         (\bm a-\bm\delta)'\bm G_H(\bm a-\bm\delta)\}.
\]
The minimizer is
\[
    \bm a_h=(\bm G_A+\bm G_H)^{-1}\bm G_H\bm\delta,
\]
so
\[
    L_h
    =\bm\delta'\bm G_H\bm\delta
      -\bm\delta'\bm G_H
       (\bm G_A+\bm G_H)^{-1}
       \bm G_H\bm\delta.
\]
Since \(\bm G_A\succeq cT\bm I_K\) and
\(\lambda_{\max}(\bm G_H)\leq h\zeta_K^2\),
\[
    \bm\delta'\bm G_H
       (\bm G_A+\bm G_H)^{-1}
       \bm G_H\bm\delta
    \leq C\frac{h\zeta_K^2}{T}
          \bm\delta'\bm G_H\bm\delta.
\]
Thus, whenever \(h\zeta_K^2/T=o(1)\),
\[
    L_h\geq\{1-o(1)\}
       \sum_{t\in H}
       \{\bm P_K(W_{t-1})'\bm\delta\}^2.
\]
A local empirical jump-energy bound
\[
    h^{-1}\sum_{t\in H}
       \{\bm P_K(W_{t-1})'\bm\delta\}^2
       \geq c\Delta_T^2
\]
therefore yields the desired \(ch\Delta_T^2\) drift.  Approximation error and
far-away candidates are precisely what the remainder
\(\rho_T(k)\) in Assumption \ref{ass:sol-localization} controls.
\end{proof}


\begin{lemma}[RSS stochastic fluctuation target]
\label{lem:wb-rss-fluctuation}
Under Assumption \ref{ass:sol-localization}, write
\[
    \Delta RSS_K(k)=RSS_K(\beta_0,k)-RSS_K(\beta_0,k_0).
\]
Uniformly over minimum-spaced one-break candidates,
\begin{align*}
    \left|\Delta RSS_K(k)
    -\E\{\Delta RSS_K(k)\mid \mathcal D_T\}\right|
    &=O_p\!\left(\sqrt{|k-k_0|K\log T}\right) \\
    &\quad +O_p(K\log T).
\end{align*}
\end{lemma}


\begin{proof}[Detailed solution]
Continue to write \(\bm A_k=\bm M_{K,k}-\bm M_{K,k_0}\) and
\(h=|k-k_0|\).  From the expansion in the preceding proof,
\begin{align}
    \Delta RSS_K(k)
    -\E\{\Delta RSS_K(k)\mid\mathcal D_T\}
    &=2\bm\mu'\bm A_k\bm U\\
    &\quad+
      \bm U'\bm A_k\bm U
      -\E(\bm U'\bm A_k\bm U\mid\mathcal D_T).
    \label{eq:sol-rss-fluctuation-decomp}
\end{align}
Conditional on \(\mathcal D_T\), the first term is sub-Gaussian with variance
proxy bounded by
\[
    C\|\bm A_k\bm\mu\|^2\leq ChK
\]
by \eqref{eq:sol-mean-sensitivity}.  Hence, for a sufficiently large constant
\(C_1\),
\[
    \Pp\left(
      |2\bm\mu'\bm A_k\bm U|>
      C_1\sqrt{hK\log T}
      \ \middle|\ \mathcal D_T
    \right)\leq 2T^{-cC_1^2}.
\]
There are at most \(T\) candidate dates, so a union bound gives, uniformly in
\(k\),
\[
    |2\bm\mu'\bm A_k\bm U|
    =O_p\{\sqrt{hK\log T}\}.
\]

For the quadratic term, \(\bm A_k\) is symmetric and
\[
    \operatorname{rank}(\bm A_k)
    \leq\operatorname{rank}(\bm\Pi_{K,k})+
       \operatorname{rank}(\bm\Pi_{K,k_0})
    \leq4K.
\]
Also \(\|\bm A_k\|\leq2\) and
\(\|\bm A_k\|_F\leq4\sqrt K\).  The conditional Hanson--Wright inequality
therefore gives, for every \(x>0\),
\[
    \Pp\left(
      |\bm U'\bm A_k\bm U-
        \E(\bm U'\bm A_k\bm U\mid\mathcal D_T)|
      >C\{\sqrt{Kx}+x\}
      \ \middle|\ \mathcal D_T
    \right)
    \leq2e^{-cx}.
\]
Taking \(x=C_2\log T\) and applying another union bound over all dates gives
\[
    \sup_k
    |\bm U'\bm A_k\bm U-
      \E(\bm U'\bm A_k\bm U\mid\mathcal D_T)|
    =O_p\{\sqrt{K\log T}+\log T\}
    =O_p(K\log T).
\]
Combining the linear and quadratic bounds in
\eqref{eq:sol-rss-fluctuation-decomp} proves
\[
    \left|\Delta RSS_K(k)
    -\E\{\Delta RSS_K(k)\mid\mathcal D_T\}\right|
    =O_p\{\sqrt{hK\log T}\}+O_p(K\log T)
\]
uniformly over the candidate set.  This logarithmic conclusion uses the
sub-Gaussian completion; it cannot in general be recovered from only the
finite-moment assumption in Section 2.
\end{proof}


\begin{theorem}[One-break localization theorem target]
\label{thm:wb-localization}
If Lemmas \ref{lem:wb-rss-drift} and \ref{lem:wb-rss-fluctuation} hold and the break signal dominates the search fluctuation, then
\[
    |\widetilde k-k_0|=O_p(r_T),
\]
where \(r_T\) is any sequence satisfying the resulting drift-fluctuation inequality. A sufficient fixed-break target is \(r_T=O_p(K\log T/\Delta_T^2)\), subject to refinement after the fluctuation bound is proved.
\end{theorem}


\begin{proof}[Detailed solution]
Let
\[
    R_T=K\log T,
    \qquad
    r_T=R_T\max\{\Delta_T^{-2},\Delta_T^{-4}\}.
    \label{eq:sol-localization-rate}
\]
The sequence \(r_T\) is deterministic; the statement
``\(r_T=O_p(\cdot)\)'' in the original target should read
``\(r_T=O(\cdot)\).''  Lemmas \ref{lem:wb-rss-drift} and
\ref{lem:wb-rss-fluctuation} imply that, with a tight random constant
\(C_T=O_p(1)\),
\[
    \Delta RSS_K(k)
    \geq c h\Delta_T^2
      -C_T\{\sqrt{hR_T}+R_T\}
      -o_p(h\Delta_T^2),
    \qquad h=|k-k_0|.
\]
Fix \(M>0\) and suppose \(h\geq Mr_T\).  By the definition of \(r_T\),
\[
    \frac{R_T}{h\Delta_T^2}\leq\frac1M
\]
and
\[
    \frac{\sqrt{hR_T}}{h\Delta_T^2}
    =\frac{\sqrt{R_T}}{\sqrt h\,\Delta_T^2}
    \leq\frac1{\sqrt M}.
\]
Choose a deterministic \(C\) so that
\(\Pp(C_T>C)<\varepsilon\), and then choose \(M\) so large that
\(C(M^{-1/2}+M^{-1})<c/4\).  On the resulting event, uniformly for
\(h\geq Mr_T\),
\[
    \Delta RSS_K(k)>0
\]
for all sufficiently large \(T\).  But the minimizer satisfies
\[
    \Delta RSS_K(\widetilde k)
    =RSS_K(\beta_0,\widetilde k)-RSS_K(\beta_0,k_0)\leq0,
\]
because \(k_0\) is an admissible candidate.  Therefore
\[
    \Pp\{|\widetilde k-k_0|>Mr_T\}\leq\varepsilon+o(1).
\]
Letting first \(T\to\infty\), then \(M\to\infty\), proves
\(|\widetilde k-k_0|=O_p(r_T)\).

If \(\Delta_T\) is bounded away from zero, then
\(r_T\asymp K\log T\), equivalently
\(r_T=O(K\log T/\Delta_T^2)\).  If \(\Delta_T\to0\), the fluctuation bound
as stated forces the additional \(\Delta_T^{-4}\) term.  The smaller rate
\(K\log T/\Delta_T^2\) for shrinking jumps requires a refined linear
fluctuation of order
\(\Delta_T\sqrt{hK\log T}\).  Finally, the score bridge requires the separate
check
\[
    K\log T\max\{\Delta_T^{-2},\Delta_T^{-4}\}=o(\sqrt T).
\]
\end{proof}


\section{Unknown-Order Selection}

\begin{assumption}[Penalty-rate target]
\label{ass:wb-penalty}
Let \(R_T\) denote the maximal stochastic RSS gain from searching over
spurious breaks; under Lemma \ref{lem:wb-rss-fluctuation}, one may take
\(R_T=K\log T\).  The penalty satisfies
\[
    \frac{\operatorname{pen}_T(q,K)-\operatorname{pen}_T(q_0,K)}{R_T}
    \to\infty
    \quad\text{for }q>q_0,
\]
and is small in absolute value relative to the deterministic RSS drift from
omitting true breaks:
\[
    \left|\operatorname{pen}_T(q,K)-\operatorname{pen}_T(q_0,K)\right|
    =o(T\Delta_T^2)
    \quad\text{for }q<q_0.
\]
These are the rates actually used in Theorem \ref{thm:wb-order}.
\end{assumption}

\begin{theorem}[Unknown-order consistency target]
\label{thm:wb-order}
Under the localization, RSS drift, RSS fluctuation, and penalty-rate targets,
\[
    \Pp\{\widehat q(\beta_0)=q_0\}\to1.
\]
\end{theorem}


\begin{proof}[Detailed solution]
Let
\[
    RSS_q^*=RSS_K\{\beta_0,\widehat\Lambda_q(\beta_0)\},
    \qquad R_T=K\log T.
\]
The RSS drift and search bounds needed for order selection can be written
uniformly over the fixed set \(0\leq q\leq\bar q\) as
\begin{align}
    \min_{q<q_0}(RSS_q^*-RSS_{q_0}^*)
      &\geq cT\Delta_T^2-O_p(R_T),
      \label{eq:sol-underfit-gap}\\
    \max_{q>q_0}(RSS_{q_0}^*-RSS_q^*)
      &=O_p(R_T).
      \label{eq:sol-overfit-gain}
\end{align}
The first display says that at least one true jump is omitted; the resulting
deterministic loss is of order \(T\Delta_T^2\).  The second says that an
extra, spurious split can exploit noise only at the searched stochastic order.

The penalty condition in the original workbook must be sharpened as follows:
for every \(q>q_0\),
\[
    \frac{\operatorname{pen}_T(q,K)-
          \operatorname{pen}_T(q_0,K)}{R_T}\longrightarrow\infty,
\]
and for every \(q<q_0\),
\[
    \left|\operatorname{pen}_T(q,K)-
          \operatorname{pen}_T(q_0,K)\right|
    =o(T\Delta_T^2).
\]
Mere divergence of the overfitting penalty is not enough when
\(R_T\to\infty\).

For \(q<q_0\), subtract the true-order penalized criterion:
\begin{align*}
    &RSS_q^*+\operatorname{pen}_T(q,K)
      -RSS_{q_0}^*-\operatorname{pen}_T(q_0,K)\\
    &\qquad\geq
      cT\Delta_T^2-O_p(R_T)-o(T\Delta_T^2)>0
\end{align*}
with probability tending to one, because
\(R_T=o(T\Delta_T^2)\).  For \(q>q_0\),
\begin{align*}
    &RSS_q^*+\operatorname{pen}_T(q,K)
      -RSS_{q_0}^*-\operatorname{pen}_T(q_0,K)\\
    &\qquad\geq
      -O_p(R_T)+
      \{\operatorname{pen}_T(q,K)-
        \operatorname{pen}_T(q_0,K)\}>0
\end{align*}
with probability tending to one.  Since \(\bar q\) is fixed, a finite union
bound proves
\[
    \Pp\{\widehat q(\beta_0)=q_0\}\to1.
\]
A transparent sufficient penalty is
\[
    \operatorname{pen}_T(q,K)=q b_T,
    \qquad
    K\log T\ll b_T\ll T\Delta_T^2.
\]
Such a sequence exists because break detectability gives
\(K\log T=o(T\Delta_T^2)\); for example,
\(b_T=\sqrt{(K\log T)(T\Delta_T^2)}\) satisfies the two rate inequalities.
\end{proof}


\begin{corollary}[Unknown-order Wilks]
\label{cor:wb-unknown-wilks}
Suppose Theorem \ref{thm:wb-order} and the fixed-order
estimated-partition Wilks theorem hold.  Then
\[
    \ell_T\{\beta_0,\widehat\Lambda(\beta_0)\}\Rightarrow\chi_1^2.
\]
\end{corollary}


\begin{proof}[Detailed solution]
Let
\[
    E_T=\{\widehat q(\beta_0)=q_0\}.
\]
By Theorem \ref{thm:wb-order}, \(\Pp(E_T)\to1\).  On \(E_T\), the definitions
and the common deterministic tie-breaking rule give
\[
    \widehat\Lambda(\beta_0)
    =\widehat\Lambda_{q_0}(\beta_0)
    =\widetilde\Lambda.
\]
For every continuity point \(x\) of the \(\chi_1^2\) distribution,
\begin{align*}
    &\left|
      \Pp\left[
        \ell_T\{\beta_0,\widehat\Lambda(\beta_0)\}\leq x
      \right]
      -\Pp\left[
        \ell_T(\beta_0,\widetilde\Lambda)\leq x
      \right]
    \right|\\
    &\qquad\leq \Pp(E_T^c)\to0.
\end{align*}
The fixed-order estimated-partition Wilks theorem then yields the claimed
\(\chi_1^2\) limit.
\end{proof}


\section{Local Alternatives}

\begin{assumption}[Local-law oracle and drift]
\label{ass:wb-local-oracle}
Let \(\beta_T=\beta_0+d_T\), \(d_T\to0\), and define
\[
    g_{t,T}=d_T(\bm M_{K,\Lambda_0}\bm X_{lag})_tv_{t,0}.
\]
Under the local law \(P_{\beta_T}\),
\[
    T^{-1/2}\sum_{t=1}^TU_tv_{t,0}\xrightarrow{\mathrm{st}}\eta_qZ,
    \qquad
    T^{-1}\sum_{t=1}^TU_t^2v_{t,0}^2\xrightarrow{p}\eta_q^2,
    \qquad
    \max_t|U_tv_{t,0}|=o_p(\sqrt T),
\]
and
\[
    T^{-1/2}\sum_{t=1}^Tg_{t,T}\xrightarrow{p}\eta_q\delta_h,
    \qquad
    T^{-1}\sum_{t=1}^Tg_{t,T}^2=o_p(1),
    \qquad
    \max_t|g_{t,T}|=o_p(\sqrt T).
\]
\end{assumption}

\begin{lemma}[Local shifted-score replacement target]
\label{lem:wb-local-replacement}
Suppose Assumptions \ref{ass:prim-sieve}, \ref{ass:prim-bounded},
\ref{ass:prim-gram}, and \ref{ass:sol-oracle-envelope} continue to hold under
\(P_{\beta_T}\).  Then
\begin{align}
    T^{-1/2}\sum_{t=1}^T
    \{\widehat Z_t(\beta_0,\Lambda_0)-(U_tv_{t,0}+g_{t,T})\}
    &=o_p(1),
    \\
    T^{-1}\sum_{t=1}^T
    \{\widehat Z_t(\beta_0,\Lambda_0)^2-(U_tv_{t,0}+g_{t,T})^2\}
    &=o_p(1),
    \\
    \max_t|\widehat Z_t(\beta_0,\Lambda_0)-(U_tv_{t,0}+g_{t,T})|
    &=o_p(\sqrt T).
\end{align}
\end{lemma}


\begin{proof}[Detailed solution]
Under \(P_{\beta_T}\),
\[
    \bm Y-\beta_0\bm X_{lag}
    =\bm Q_{K,\Lambda_0}\bm c_0+
      \bm r_0+\bm U+d_T\bm X_{lag}.
\]
Apply \(\bm M_{K,\Lambda_0}\) and put
\[
    \bm e_0=\bm M_{K,\Lambda_0}\bm r_0,
    \qquad
    \bm h_0=\bm\Pi_{K,\Lambda_0}\bm U,
    \qquad
    \bm x_0^c=\bm M_{K,\Lambda_0}\bm X_{lag}.
\]
Then
\[
    \widehat{\bm u}(\beta_0,\Lambda_0)
    =\bm U-\bm h_0+\bm e_0+d_T\bm x_0^c.
\]
Multiplying componentwise by \(\bm v_0\) gives the exact identity
\[
    \widehat Z_t(\beta_0,\Lambda_0)
    =U_tv_{t,0}+g_{t,T}+(e_{t,0}-h_{t,0})v_{t,0}.
\]
The final term is exactly the same feasibility error as under the null.  The
proof of Lemma \ref{lem:wb-feasible-replacement} therefore gives
\[
    \|(\bm e_0-\bm h_0)\odot\bm v_0\|_T=o_p(1),
    \qquad
    \|(\bm e_0-\bm h_0)\odot\bm v_0\|_\infty=o_p(\sqrt T),
\]
and the numerator projection cancellation plus the projection-bias condition
gives its normalized sum as \(o_p(1)\).  This proves the first and third
displays of the lemma.

For the square replacement, put
\(a_{t,T}=U_tv_{t,0}+g_{t,T}\) and
\(d_{t,T}=(e_{t,0}-h_{t,0})v_{t,0}\).  Then
\[
    T^{-1}\sum_t
      \{\widehat Z_t^2-a_{t,T}^2\}
    =2T^{-1}\sum_ta_{t,T}d_{t,T}
      +T^{-1}\sum_td_{t,T}^2.
\]
Assumption \ref{ass:wb-local-oracle} gives
\[
    \|\bm U\odot\bm v_0\|_T=O_p(1),
    \qquad
    \|\bm g_T\|_T=o_p(1),
\]
so \(\|\bm a_T\|_T=O_p(1)\).  Cauchy--Schwarz and
\(\|\bm d_T\|_T=o_p(1)\) make both terms on the right \(o_p(1)\), proving the
second display.
\end{proof}


\begin{theorem}[Local noncentral chi-square target]
\label{thm:wb-local-power}
Under Assumption \ref{ass:wb-local-oracle}, Lemma \ref{lem:wb-local-replacement}, the local convex-hull condition, and the estimated-partition bridge under \(P_{\beta_T}\),
\[
    \ell_T(\beta_0,\widetilde\Lambda)
    \Rightarrow
    \chi_1^2(\delta_h^2).
\]
If \(\Pp\{\widehat q(\beta_0)=q_0\}\to1\), the same limit holds for \(\ell_T\{\beta_0,\widehat\Lambda(\beta_0)\}\).
\end{theorem}


\begin{proof}[Detailed solution]
First work at the true partition.  Lemma \ref{lem:wb-local-replacement} and
Assumption \ref{ass:wb-local-oracle} give
\begin{align*}
    T^{-1/2}\sum_t\widehat Z_t(\beta_0,\Lambda_0)
    &=T^{-1/2}\sum_tU_tv_{t,0}
      +T^{-1/2}\sum_tg_{t,T}+o_p(1)\\
    &\xrightarrow{\mathrm{st}}\eta_q(Z+\delta_h).
\end{align*}
For the second moment,
\begin{align*}
    T^{-1}\sum_t(U_tv_{t,0}+g_{t,T})^2
    &=T^{-1}\sum_tU_t^2v_{t,0}^2\\
    &\quad+2T^{-1}\sum_tU_tv_{t,0}g_{t,T}
      +T^{-1}\sum_tg_{t,T}^2.
\end{align*}
The first term converges to \(\eta_q^2\), and the last is \(o_p(1)\).  By
Cauchy--Schwarz,
\[
    \left|T^{-1}\sum_tU_tv_{t,0}g_{t,T}\right|
    \leq
    \left(T^{-1}\sum_tU_t^2v_{t,0}^2\right)^{1/2}
    \left(T^{-1}\sum_tg_{t,T}^2\right)^{1/2}
    =o_p(1).
\]
The local replacement lemma therefore yields
\[
    T^{-1}\sum_t\widehat Z_t(\beta_0,\Lambda_0)^2
    \xrightarrow{p}\eta_q^2.
\]
Moreover,
\[
    \max_t|U_tv_{t,0}+g_{t,T}|=o_p(\sqrt T),
\]
and the maximum replacement error is also \(o_p(\sqrt T)\).  The local
convex-hull condition and Lemma \ref{lem:sol-scalar-el} imply
\[
    \ell_T(\beta_0,\Lambda_0)
    \Rightarrow
    \frac{\eta_q^2(Z+\delta_h)^2}{\eta_q^2}
    =(Z+\delta_h)^2
    \sim\chi_1^2(\delta_h^2).
\]

Under the local-law version of the one-break bridge, the estimated-partition
numerator and denominator differ from their true-partition counterparts by
\(o_p(1)\), and the estimated score maximum is \(o_p(\sqrt T)\).  Applying the
same scalar EL expansion transfers the limit to
\(\ell_T(\beta_0,\widetilde\Lambda)\).  Finally, on the event
\(\{\widehat q(\beta_0)=q_0\}\), the unknown-order and fixed-order partitions
coincide.  If that event has probability tending to one under
\(P_{\beta_T}\), the same event-transfer argument as in Corollary
\ref{cor:wb-unknown-wilks} proves the final assertion.
\end{proof}


\section{Master Dependency Checklist}

\begin{center}
\begin{tabular}{>{\raggedright\arraybackslash}p{0.28\textwidth}
                >{\raggedright\arraybackslash}p{0.34\textwidth}
                >{\raggedright\arraybackslash}p{0.28\textwidth}}
\hline
Goal & Must prove first & Output when finished \\
\hline
Known-partition Wilks & Lemmas \ref{lem:wb-oracle-clt}, \ref{lem:wb-oracle-lln}, \ref{lem:wb-oracle-max}, \ref{lem:wb-feasible-replacement} & Theorem \ref{thm:wb-known-wilks} \\
One-break bridge & Lemmas \ref{lem:wb-bridge-mu}, \ref{lem:wb-bridge-u}, \ref{lem:wb-bridge-v} & Theorem \ref{thm:wb-one-break-equivalence} \\
Break localization & Lemmas \ref{lem:wb-rss-drift}, \ref{lem:wb-rss-fluctuation} & Theorem \ref{thm:wb-localization} \\
Unknown order & Theorem \ref{thm:wb-order} and penalty-rate target & Corollary \ref{cor:wb-unknown-wilks} \\
Local alternatives & Local oracle, shifted replacement, local bridge & Theorem \ref{thm:wb-local-power} \\
\hline
\end{tabular}
\end{center}

\subsection*{Recommended homework order}
\begin{enumerate}
    \item Prove deterministic projection algebra and leverage bounds.
    \item Prove known-partition Wilks from primitive martingale assumptions.
    \item Prove the one-break partition perturbation bounds.
    \item Prove profile-RSS break localization.
    \item Only then work on unknown-order selection and local alternatives.
\end{enumerate}

\end{document}
